\partnote{Resolution, repair, and retained detail}{Resolution gives a concrete family of views. Its compatibility differential connects exact maps, protected information, and residual coordinates. Recovering every finite coordinate is then separated from convergence in a chosen norm.}
\section{Resolution towers and exact transport}\label{sec:towers}
\subsection{The sequential model}
A tower in a category \(\mathcal C\) is a sequence
\[
 \cdots\xrightarrow{\rho_2}X_2\xrightarrow{\rho_1}X_1
 \xrightarrow{\rho_0}X_0.
\]
For sets, a point is a sequence \(x_n\in X_n\) with \(\rho_n(x_{n+1})=x_n\). Composite restrictions turn the tower into a view system on \(\N\), with an arrow \(i\to j\) when \(i\le j\). Its section space is the same coherent point type, with the same coordinates.

A candidate transport from \(X\) to \(Y\) consists of maps \(f_n:X_n\to Y_n\). It is exact when
\begin{equation}\label{eq:exacttower}
 \rho^Y_n f_{n+1}=f_n\rho^X_n\quad\text{for all }n.
\end{equation}

\begin{proposition}[The exact tower category]\label{KS16}
Identity maps and composites of exact transports are exact. An exact transport induces a map of coherent point spaces by \(x\mapsto(f_nx_n)_n\), preserving identities and composition.
\end{proposition}
\begin{proof}
For two exact transports, substitute \(\rho^Z_ng_{n+1}=g_n\rho^Y_n\) and \(\rho^Y_nf_{n+1}=f_n\rho^X_n\) into the composite square. For a coherent point,
\[
 \rho^Y_n f_{n+1}x_{n+1}
 =f_n\rho^X_nx_{n+1}=f_nx_n.
\]
Identity, associativity, and composition on points are checked coordinatewise.
\end{proof}
\source{KS16}
This is a category on a fixed schedule. A full procategory allows additional reindexing and equivalence relations on maps. Those constructions must be supplied rather than identified with the fixed schedule by terminology.

\subsection{The actual residual of a transport}
Assume \(\mathcal C\) is preadditive, so differences of maps make sense. The tower differential is
\[
 (\partial f)_n=\rho^Y_nf_{n+1}-f_n\rho^X_n.
\]
Its target has components \(\Hom(X_{n+1},Y_n)\), not \(\Hom(X_n,Y_n)\).

\begin{proposition}[The transport Leibniz identity]\label{KS17}
For composable candidates,
\[
 \partial(gf)_n=(\partial g)_nf_{n+1}+g_n(\partial f)_n.
\]
In particular exactness is equivalent to vanishing of the differential.
\end{proposition}
\begin{proof}
Expand the right side. The terms \(-g_n\rho^Y_nf_{n+1}\) and \(+g_n\rho^Y_nf_{n+1}\) cancel, leaving \(\rho^Z_ng_{n+1}f_{n+1}-g_nf_n\rho^X_n\). This is the differential of the composite. The exactness statement is \eqref{eq:exacttower}.
\end{proof}
\source{KS17}
There is no additional mixed defect term in this identity. Such a term would require a different definition of the defect or a different expansion. Typed cancellation is a useful safeguard against importing an attractive but incorrect error formula.

\subsection{Cofinal views preserve points, not every measurement}
Let \(\phi:\N\to\N\) be monotone and unbounded. The sampled tower uses levels \(X_{\phi(n)}\) and their composite restrictions.

\begin{theorem}[Cofinal reconstruction of points]\label{KS18}
Sampling induces an equivalence between the coherent points of \(X\) and of its sampled tower.
\end{theorem}
\begin{proof}
For a coherent sampled point \(y\) and an arbitrary level \(k\), choose \(n\) with \(k\le\phi(n)\) and restrict \(y_n\) to \(X_k\). Two choices agree after passing to a later sampled stage, using monotonicity and coherence. These recovered coordinates satisfy every bonding equation. Restricting a full point and recovering it gives its original coordinates, while sampling the recovered family gives \(y\).
\end{proof}
\source{KS18}
This theorem is about exact coherent points. Raw residual size is a different quantity. On identity scalar towers, \(f_n=(-1)^n\) has nonzero fine differences, but its even subsequence is constant. The coarse residual vanishes because adjacent fine defects cancel. A cofinal equivalence of limits does not imply equality of residual norms or an absence of aliasing in uncorrected data.

\plate{03_tower}{Exact transport commutes with every adjacent restriction. A cofinal subsequence retains coherent points, but a raw defect can cancel during coarsening. The retained coarse value alone need not record the omitted fine discrepancy.}{fig:tower}

\chap{Banach repair and constructed split coordinates}\label{sec:split}
\subsection{An actual complete candidate space}
Let \(X_n,Y_n\) be Banach spaces over a complete nontrivially normed field, with continuous linear bonding maps and uniform bounds \(\norm{\rho^X_n}\le B_X\), \(\norm{\rho^Y_n}\le B_Y\). Use the spaces of bounded operator families
\[
 E=\ell^\infty_n\mathcal L(X_n,Y_n),\qquad
 R=\ell^\infty_n\mathcal L(X_{n+1},Y_n).
\]
The source formalization derives the required completeness from the operator spaces and their bounded product. The same differential satisfies
\[
 \norm{\partial f}_\infty\le(B_X+B_Y)\norm f_\infty.
\]
Thus it is a continuous linear map, and evaluation at a fixed level is continuous. These facts connect algebraic compatibility to an analytic limit argument.

\subsection{From a correction operator to a projection}
The next statement is general. Let \(E\) be Banach, let \(d:E\to R\), \(p:E\to S\), and \(A:R\to E\) be continuous linear maps. Suppose \(pA=0\) and, for \(0\le q<1\),
\begin{equation}\label{eq:contraction}
 \norm{df-dA(df)}\le q\norm{df}\qquad(f\in E).
\end{equation}
The hypothesis concerns residuals of actual candidates. It need not hold for every independent vector in \(R\).

\begin{theorem}[Constructed geometric repair]\label{KS19}
Starting from \(f^{(0)}=f\), define
\[
 f^{(m+1)}=f^{(m)}-A(df^{(m)}).
\]
There is a bounded linear projection \(P_A:E\to\ker d\) with
\[
 f^{(m)}\longrightarrow P_Af,\qquad p(P_Af)=p(f),
\]
and
\[
 \norm{f^{(m)}-P_Af}
 \le\frac{\norm A}{1-q}\norm{df^{(m)}}.
\]
It fixes every exact candidate and satisfies
\[
 \norm{P_A}\le1+\frac{\norm A\norm d}{1-q}.
\]
\end{theorem}
\begin{proof}
Equation~\eqref{eq:contraction} gives geometric decay of the residuals. The step size is at most \(\norm A\norm{df^{(m)}}\), so the series of movements converges and the sequence is Cauchy. Completeness supplies a limit. Continuity of \(d\) gives zero residual, and \(pA=0\) makes every iterate preserve \(p(f)\). Summing the tail of the geometric movement estimate gives the displayed error bound.

Each iterate is linear in its initial value. Addition and scalar multiplication commute with the limit, so the limit map is linear. The estimate at \(m=0\), followed by the triangle inequality, gives its operator bound. If \(df=0\), the iteration is stationary. Therefore its image is exactly \(\ker d\), it fixes that image, and it is idempotent.
\end{proof}
\source{KS19}
This theorem constructs the dynamics from \(A\). It does not construct a suitable \(A\) for every residual problem. That missing input can, however, be derived for a useful class of actual towers.

\subsection{A right inverse from target splittings}
Suppose each target bonding map has a continuous linear section
\[
 s_n:Y_n\to Y_{n+1},\qquad \rho^Y_ns_n=I,\qquad \norm{s_n}\le S.
\]
Assume the propagation gap
\begin{equation}\label{eq:splitgap}
 SB_X<1,\qquad K=\frac{S}{1-SB_X}.
\end{equation}
For \(r\in R\), recursively set
\begin{equation}\label{eq:greenrecurrence}
 (Ar)_0=0,\qquad
 (Ar)_{n+1}=s_n\bigl(r_n+(Ar)_n\rho^X_n\bigr).
\end{equation}

\begin{theorem}[Split tower correction]\label{KS20}
The recurrence defines a continuous linear map \(A:R\to E\) with
\[
 \partial A=I_R,\qquad \norm A\le K.
\]
The operator \(P=I-A\partial\) is an exact projection and
\[
 (Pf)_0=f_0,\qquad
 \norm{f-Pf}_\infty\le K\norm{\partial f}_\infty.
\]
\end{theorem}
\begin{proof}
Inductively, if \(\norm{(Ar)_n}\le K\norm r_\infty\), then
\[
 \norm{(Ar)_{n+1}}
 \le S(1+B_XK)\norm r_\infty=K\norm r_\infty.
\]
The base case is zero. The recurrence is linear, so this uniform bound bundles it into a bounded linear operator. Applying \(\rho^Y_n\) to \eqref{eq:greenrecurrence} and using the section identity gives
\[
 (\partial Ar)_n=r_n+(Ar)_n\rho^X_n-(Ar)_n\rho^X_n=r_n.
\]
Consequently \(\partial P=0\), \(P^2=P\), and the error is \(A\partial f\). The initial correction vanishes.
\end{proof}
\source{KS20}
The surjectivity is onto all bounded residual families. It is stronger than the assertion that some exact output can be assigned to each candidate. The gap and the chosen splittings are genuine inputs.

\subsection{Protected information beyond the seed}
When the target grows, its new kernel directions are independent information. Define
\[
 N_n(f)=(I-s_n\rho^Y_n)f_{n+1},\qquad
 (p_Lf)_0=f_0,\quad(p_Lf)_{n+1}=N_n(f).
\]
Each \(N_n(f)\) lands in \(\ker\rho^Y_n\). The observable \(p_L\) retains both the seed and every new output component, relative to the chosen sections.

\begin{theorem}[Split coordinate theorem]\label{KS21}
The observable is a bounded projection with \(p_LA=0\) and \(p_LP=p_L\). It separates exact candidates. If
\[
 \mathcal S_L=\{s\in E:p_Ls=s\},
\]
then
\[
 \Phi:E\simeq_L\mathcal S_L\times R,
 \qquad \Phi(f)=(p_Lf,\partial f),
 \qquad \Phi^{-1}(s,r)=Ps+Ar.
\]
Repair becomes \((s,r)\mapsto(s,0)\). It is the unique exact repair preserving \(p_Lf\).
\end{theorem}
\begin{proof}
The operator \(I-s_n\rho^Y_n\) is a projection onto the target kernel. This proves idempotence of \(p_L\), and its norm is at most \(1+SB_Y\). Every correction in \eqref{eq:greenrecurrence} at a positive level lies in the range of the corresponding section, so its innovation is zero. Hence \(p_LA=0\) and \(p_LP=p_L\).

If \(f\) is exact, then
\[
 f_{n+1}=s_nf_n\rho^X_n+N_n(f).
\]
Its seed and innovations determine the family by induction. Therefore two exact families with the same observable agree. More generally, if \(f,g\) have the same observable and residual, apply this fact to the exact difference \(f-g\).

For \((s,r)\in\mathcal S_L\times R\), the candidate \(Ps+Ar\) has observable \(s\) and residual \(r\). Conversely \(f\) and \(P(p_Lf)+A\partial f\) have the same two coordinates, so they agree. Boundedness of \(p_L\), \(\partial\), \(P\), and \(A\) makes both inverse maps continuous. Finally, zeroing the residual retains \(s\), and separation proves uniqueness.
\end{proof}
\source{KS21}
This is an actual analytic coordinate equivalence, not generally an isometry. The protected coordinate depends on the declared splitting. It should not be called an intrinsic semantic coordinate without a further comparison theorem.

\plate{04_split}{In the split coordinate chart, a candidate separates into retained information and compatibility residual. The exact subspace is the zero residual slice. The chart and its protected coordinate depend on the chosen target sections.}{fig:split}

\subsection{A growing example and the gap boundary}
Use \(X_n=Y_n=\R^{n+1}\) with maximum norms. Let \(\pi_n\) delete the last coordinate, take \(\rho^X_n=\tfrac12\pi_n\), \(\rho^Y_n=\pi_n\), and let \(s_n\) append zero. Then \(S=B_Y=1\), \(B_X=1/2\), and \(K=2\).

\begin{example}[Repair of the identity family]\label{KS22}
The bounded candidate \(f_n=I\) is not exact. Its repair is
\[
 (Pf)_n=\operatorname{diag}(2^{-n},2^{-(n-1)},\ldots,2^{-1},1).
\]
The new last output row is retained at every stage while inherited rows are propagated by the factor \(1/2\).
\end{example}
\begin{proof}
At level zero the value is one. Exactness forces every inherited row at the next level to be one half of its previous value. Preservation of the innovation forces the new row to remain the corresponding identity row. Induction gives the displayed diagonal.
\end{proof}
\source{KS22}
For constant identity towers, by contrast, \((\partial f)_n=f_{n+1}-f_n\). The constant residual \(r_n=1\) has no bounded primitive, since any primitive satisfies \(f_n=f_0+n\). Here \(SB_X=1\). This counterexample excludes discarding the gap indiscriminately. It does not prove that the gap is necessary for every individual tower. Exact repair preserving the initial coordinate still exists in this constant example, by replacing a candidate with its initial constant family. Surjectivity onto all bounded residuals does not follow from that repair.
\source{KS23}

\subsection{What is canonical}
For fixed splitting and protected observable, the repair above is unique. Changing the splitting can change both the observable and the repaired answer. In the growing example, replace \(s_n(y)=(y,0)\) by \(s_n'(y)=(y,y_0)\). At the first positive level, repair of the identity becomes
\[
 \begin{pmatrix}1/2&0\\-1/2&1\end{pmatrix}
 \quad\text{instead of}\quad
 \begin{pmatrix}1/2&0\\0&1\end{pmatrix}.
\]
Both are exact continuations for their respective innovation profiles and have the same seed. Thus a unique answer for a declared preservation problem is not an answer independent of all preservation choices.
\source{KS24}

\chap{Multiscale residuals and the limits of reconstruction}\label{sec:cascade}
\subsection{A split coarsener retains its missing detail}
Let \(C:R\to\overline R\) and \(J:\overline R\to R\) be linear maps with \(CJ=I\). Put \(H=I-JC\). Then
\[
 r=JCr+Hr,\qquad CH=0,\qquad H^2=H.
\]
The two coordinates \((Cr,Hr)\) retain all information. The coarse coordinate alone generally does not.

The actual dyadic tower coarsener has the typed formula
\begin{equation}\label{eq:dyadic}
 (Cr)_n=\rho^Y_{2n}r_{2n+1}+r_{2n}\rho^X_{2n+1}.
\end{equation}
It is obtained by telescoping two adjacent compatibility squares. With target sections, a lift is
\[
 (Jc)_{2n}=0,\qquad (Jc)_{2n+1}=s_{2n}c_n.
\]
The section identity gives \(CJ=I\), so the missing detail is explicitly \(Hr\), not an unspecified reconstruction error. The original normed construction imposes the bounds needed to make these maps bounded on its operator family spaces.
\source{KS25}

\subsection{A whole cascade}
At scale \(j\), let \(C_j:R_j\to R_{j+1}\) and \(J_j:R_{j+1}\to R_j\) be bounded linear maps with \(C_jJ_j=I\). For \(x\in R_0\), define
\[
 x_0=x,\quad x_{j+1}=C_jx_j,\quad
 L_0=I,\quad L_N=J_0\cdots J_{N-1},
\]
\[
 h_j=x_j-J_jx_{j+1},\qquad Q_Nx=L_Nx_N,
 \qquad S_N=I-Q_N.
\]
The map \(Q_N\) restores the surviving coarse remainder to the original space. It is this restored remainder, not just \(x_N\), that controls reconstruction.

\begin{theorem}[Finite reconstruction and nested remainders]\label{KS26}
For every \(N\),
\[
 x=Q_Nx+\sum_{j<N}L_jh_j,
 \qquad S_Nx=\sum_{j<N}L_jh_j.
\]
Each \(Q_N\) is a projection, and for \(N\le M\),
\[
 Q_NQ_M=Q_MQ_N=Q_M.
\]
The vectors with all details zero are exactly the persistent subspace
\[
 \mathcal P_\infty=\{x:Q_Nx=x\text{ for every }N\}.
\]
\end{theorem}
\begin{proof}
The one scale identity gives \(L_jh_j=Q_jx-Q_{j+1}x\). Summing telescopes from \(Q_0=I\). A coarsening composite followed by its matching restoration cancels because \(C_jJ_j=I\). This proves idempotence and both nested projection identities. If all details vanish, successive remainders are equal and hence all equal \(x\). The converse follows from their successive differences.
\end{proof}
\source{KS26}

\plate{05_cascade}{A split cascade retains every discarded component as a detail. The finite sum is always exact after the restored coarse remainder is included. An infinite detail expansion requires that remainder to vanish in the topology being used.}{fig:cascade}

\subsection{Convergence has its own theorem}
Call \(x\) norm recoverable when \(Q_Nx\to0\). This is equivalent to \(S_Nx\to x\), but it does not follow merely from vanishing persistent subspace. Let
\[
 \mathcal F=\bigcup_N\ker Q_N
\]
be the finite core.

\begin{theorem}[Density criterion for norm recovery]\label{KS27}
Suppose \(\sup_N\norm{Q_N}\le K<\infty\). Then the norm recoverable vectors are exactly \(\overline{\mathcal F}\). In particular all vectors are norm recoverable if and only if the finite core is dense.
\end{theorem}
\begin{proof}
If \(Q_Nx\to0\), then \(S_Nx\to x\), and \(Q_NS_Nx=0\), so each approximant lies in the finite core. Conversely choose \(y\in\ker Q_M\) close to \(x\). For \(N\ge M\), nesting gives \(Q_Ny=0\), so
\[
 \norm{Q_Nx}\le K\norm{x-y}.
\]
Choose the approximation first and then the depth. This proves convergence for every \(x\) in the closure.
\end{proof}
\source{KS27}
There are other sufficient conditions. If the transported details are absolutely summable in the Banach space \(R_0\), their sum exists and the remaining boundary lies in \(\mathcal P_\infty\). If that persistent subspace is zero, the sum equals \(x\). Alternatively,
\[
 \norm{x-S_Nx}\le\left(\prod_{j<N}\norm{J_j}\right)\norm{x_N}
\]
gives a direct sufficient decay test. A decreasing coarse norm by itself is not enough when restoration gains grow.

\subsection{The actual dyadic result}
For the source's dyadic residual hierarchy, the persistent residual is zero. Zero detail at a scale reveals every even coordinate, while an odd coordinate is passed to a smaller index at the next scale. Strong induction on the coordinate index therefore proves that a residual with all details zero is zero.

\begin{theorem}[Dyadic determination and finite coordinate recovery]\label{KS28}
In the actual dyadic system, all detail coordinates jointly determine the residual. Moreover,
\[
 k+1<2^N\quad\Longrightarrow\quad
 (Q_Nr)_k=0,\qquad(S_Nr)_k=r_k.
\]
Thus every fixed residual coordinate is recovered exactly at a sufficiently large finite depth. The statement does not assert convergence in the supremum norm over all coordinates.
\end{theorem}
\begin{proof}[Proof outline]
At one scale, the hidden even coordinate is the original even coordinate. An odd hidden equation expresses the remaining coordinate through the next scale. Repeated descent reduces an index \(k\) to \(\lfloor k/2\rfloor\). After \(N\) steps with \(k+1<2^N\), no unrecovered coordinate remains at that index. This proves the displayed support statement, which is formalized with the dependent tower operator types. Applying it to a residual with zero detail at every scale forces every coordinate to vanish.
\end{proof}
\source{KS28}
The strict inequality matters. At the final coordinate of a dyadic block, the remainder can still be nonzero. This is a support theorem, not an argument based on finitely many numerical samples.

\subsection{Two boundary examples}
On \(\ell^\infty(\N,\R)\), let \(C\) remove the first coordinate and let \(J\) prepend zero. Then \(Q_Nx\) is the tail of \(x\), placed at its original coordinates with zeros before \(N\). Every fixed coordinate is eventually zero in the remainder, and \(\norm{Q_N}\le1\). Nevertheless the constant sequence \(\mathbf1\) satisfies \(\norm{Q_N\mathbf1}=1\) for every \(N\).

\begin{proposition}[The recovery domain is \(c_0\)]\label{KS29}
For this shift cascade,
\[
 S_Nx\to x\text{ in }\ell^\infty
 \quad\Longleftrightarrow\quad x_n\to0.
\]
For \(x_n=q^n\), \(0\le q<1\), the error is exactly \(q^N\).
\end{proposition}
\begin{proof}
The error norm is \(\sup_{n\ge N}|x_n|\). It tends to zero precisely when the sequence tends to zero. For a geometric sequence its largest tail value is the first one, \(q^N\).
\end{proof}
\source{KS29}
As a second example, take scalar spaces with \(C=\tfrac12I\), \(J=2I\). The coarse value \(x_N=2^{-N}x\) tends to zero, but restoration gives \(Q_Nx=x\). Every detail is zero and the entire input is persistent. The example belongs to this scalar cascade, not to the actual dyadic hierarchy, whose persistent subspace was proved to vanish.

\begin{insight}
Packets of retained detail are useful because they make omission explicit. They do not automatically make an infinite reconstruction convergent. Coordinate separation, bounded reconstruction, density, summability, and attainment remain different mathematical properties.
\end{insight}
